\section{Sparse canonical radix conversion}\label{sec:radix}

Let $P=X^M+h$, where $h$ has at most two monomials with coefficients
in $\F_4$, and $d=\deg h$ satisfies $d\ell\leq M$.
Write $f=\sum_{t<\ell}f_tX^{Mt}$, with $\deg f_t<M$.

\begin{lemma}[Canonical digits and synthesis]\label{lem:radix}
Define
\[
 F_s=\sum_{t\geq s}\binom ts f_t h^{t-s}
      =L_s+X^M B_s,\qquad \deg L_s<M.
\]
With $B_{-1}=B_{\ell-1}=0$, the canonical $P$-radix digits are
\begin{equation}\label{eq:radix}
 c_s=L_s+hB_s+B_{s-1},\qquad f=\sum_{s<\ell}c_sP^s,
 \qquad \deg c_s<M.
\end{equation}
Conversely, apply the same row Taylor map to $(c_s)$, obtaining
$G_t=L'_t+X^MH_t$. The ordinary monomial blocks are
\begin{equation}\label{eq:radix-inverse}
 f_t=L'_t+H_{t-1}.
\end{equation}
\end{lemma}
\begin{proof}
Substitute $X^M=P+h$ and expand by the binomial theorem. This gives
$f=\sum_sF_sP^s$. Replace $X^MB_s$ by $(P+h)B_s$ and collect
equal powers of $P$ to obtain \eqref{eq:radix}. The degree bound
\[
 \deg F_s<M+d(\ell-1-s),\qquad
 \deg B_s<d(\ell-1-s)
\]
implies $\deg(hB_s)<d(\ell-s)\leq M$.
For synthesis expand $P^s=(X^M+h)^s$, giving
$\sum_tG_tX^{Mt}$, and collect adjacent high and low blocks.
This proves \eqref{eq:radix-inverse}; there is no $hH_t$ term in it.
\end{proof}

The condition holds for both polynomials in
Lemma~\ref{lem:annihilator}: for $P_1$, $d=1$ and $\ell\leq Q=M$;
for $P_2$, $d=Q$ and $\ell\leq Q$, whereas $M=Q^2$.

For a monomial $h=\alpha X^e$, its row Taylor operation factors as
\begin{equation}\label{eq:monomial-taylor}
 (F_s)_s=
 \operatorname{diag}(\alpha^{-s}X^{-es})\,
 \mathcal P_\ell\,
 \operatorname{diag}(\alpha^tX^{et})\,(f_t)_t.
\end{equation}
The negative shifts mean deletion of known zero prefixes, not Laurent
coefficients. Indeed every contributing index $t$ obeys $t\geq s$.
If $h=h_1+h_2$, compose the commuting Taylor translations for the two
monomials. Each intermediate row has $O(M)$ coefficients: even the
crude bound using both shifts is below $3M$. A common $4M$-slot
rectangle therefore suffices, with a constant volume factor that is
removed before any recursive evaluator call.

Store complete coefficient words contiguously. Row shifts, fixed
$\F_4$ scalings and the carry operation cost $O(M\ell D)$.
Lemma~\ref{lem:rectangle} reorients the row/coefficient address
fields at cost $O(M\ell D k^\gamma)$, where their total address
width is $O(k)$, not merely $O(r)$. On that rectangle the Pascal
map has payload equal to the entire coefficient row and is covered
by Contract~\ref{con:tensor}. Alternatively this is the direct
uniform-payload interface to that contract. Both layouts include
their physical rearrangements.

All transposes are explicit coordinate operations. If a row map is
written $R\mathcal P_\ell G$, its transpose is
$G^{\T}\mathcal P_\ell^{\T}R^{\T}$: $R^{\T}$ inserts zeros at the
deleted prefix and $G^{\T}$ selects the original input interval.
The carry map \eqref{eq:radix} sends a covector $z_s$ to $z_s$ on
$L_s$ and $M_h^{\T}z_s+z_{s+1}$ on $B_s$, with the appropriate
degree projections. Synthesis sends $z_{s+1}$ to $H_s$.
The transpose of multiplication by a fixed monomial $h$ is just a
reverse shift and fixed-field scalar transpose. These statements hold
for arbitrary covectors: the primal zero-prefix fact is never imposed
on a cotangent.

